\label{gu:app:splitting}
The preceding arguments use $D=(\chi,1,1)$. Here we give the twisting
formulas and the corresponding theta quotient statement for general
fixed data.

\subsection{Extensions and character twists}
Any oscillator extension implementing the same Heisenberg
automorphisms and agreeing with $\omega$ on $G_\theta^1\times (G'_\theta)^1$ differs
from $\omega_{R,0}$ by scalars, by the Stone--von Neumann theorem.
The representation law makes these scalars a character of
$R/(G_\theta^1\times (G'_\theta)^1)\simeq F^\times$. Thus
\begin{equation}\label{gu:jac:extension-twist}
 \omega_{R,\kappa}=\omega_{R,0}(\kappa\circ\lambda_R).
\end{equation}
Here $\lambda_R$ is the common multiplier. The extensions trivial
on $K_E$ are those with $\kappa|_{N_{E/F}E^\times}=1$. The matching Levi
normalization \eqref{gu:jac:normalized-levi} selects $\kappa=1$.

Let $D=(\chi_V,\chi_W,\kappa)$, with $\chi_V,\chi_W$ unitary
characters trivial on $F^\times$ and $\kappa$ smooth on $F^\times$.
Using \eqref{gu:eq:ambient-character}, put
\begin{equation}\label{gu:eq:A-C}
 A(g)=\mathcal A_{\chi_W}(g)\kappa(\lambda_{G_\theta}(g)),\qquad
 C(h)=\mathcal A_{\chi_V}(h),\qquad
 \omega_{R,D}=\omega_{R,0}(A\boxtimes C)|_R.
\end{equation}
Its ordinary restriction has splitting characters $(\chi_V,\chi_W)$.
Define $\Omega_D^{\GU}=\cInd_R^{G_\theta\times G'_\theta}\omega_{R,D}$ and
$\Theta_D^{\GU}(\tau)=(\Omega_D^{\GU}\otimes\tau^\vee)_{G_\theta}$.
Multiplication of inducing functions by the ambient character preserves
support and gives
\begin{equation}\label{gu:eq:ambient-twist}
 \Omega_D^{\GU}\simeq\Omega_0^{\GU}\otimes(A\boxtimes C),\qquad
 \Theta_D^{\GU}(\tau)
 \simeq C\otimes\Theta_0^{\GU}(\tau\otimes A^{-1}).
\end{equation}
Write $\Omega_\kappa^{\GU}=\Omega_{(1,1,\kappa)}^{\GU}$.
For the pure multiplier case $D=(1,1,\kappa)$, specialize
\eqref{gu:eq:ambient-twist} to $A=\kappa\circ\lambda_{G_\theta}$ and $C=1$.
Writing $\lambda_{H_j}$ for the common multiplier on $H_j$,
Proposition~\ref{gu:jac:jacquet} gives
\[
 (\Omega_\kappa^{\GU})_{N,\psi_B}
 \simeq\cInd_{H_j}^{G_\theta\times K}(\kappa\circ\lambda_{H_j}),
\]
with intertwiner $\kappa(\lambda(g))\mathcal I_j$ in the same coordinates.

The scalar formula $\mathcal A_\eta(zI_4)=\eta(z)^2$ gives the
character of $K_E$:
\begin{equation}\label{gu:eq:kernel-character}
 k_D(z)=\chi_W(z)^2\chi_V(z)^{-2}\kappa(N_{E/F}z).
\end{equation}
Every nonzero quotient $\tau\boxtimes\Pi$ of $\Omega_D^{\GU}$ thus has
\begin{equation}\label{gu:eq:theta-central}
 \omega_\tau(z)\omega_\Pi(\bar z)=k_D(z).
\end{equation}
For normalized data the scalar character on the nonzero big lift
is $\omega_{\Theta_0^{\GU}(\tau)}(z)=\omega_\tau(\bar z)^{-1}$,
directly from the tensor kernel.

\begin{proposition}[Shalika period identity for general splitting data]
\label{gu:thm:hom-general}
For these data and irreducible admissible $\tau$ and $\sigma$,
\begin{equation}\label{gu:eq:hom-general}
 \begin{aligned}
 &\Hom_{S_B^{\GU}}(\Theta_D^{\GU}(\tau),\sigma\boxtimes\psi_B)\\
 &\quad\simeq\Hom_{H_{\mathrm U}}\left(
   (\tau^\vee\otimes A)\boxtimes
   ((\sigma^j)^\vee\otimes(\chi_V^{-1}\circ\det_E)),\mathbf1\right).
 \end{aligned}
\end{equation}
The dimension is at most one.
\end{proposition}
\begin{proof}
On the Shalika subgroup,
$C(m_B(a))=\chi_V(\det_Ea)$ and $C|_N=1$. We apply
\eqref{gu:eq:ambient-twist} and the normalized case of
Theorem~\ref{gu:thm:hom} to $\tau\otimes A^{-1}$ and
$\sigma\otimes(\chi_V^{-1}\circ\det_X)$. The transport and duality
in \eqref{gu:eq:coefficient-three-steps}, with $\chi_V$ in place of
$\chi$, give \eqref{gu:eq:hom-general}; Theorem~\ref{thm:main-unitary} gives
the bound.
\end{proof}
Equivalently, the period character before absorbing the twists is
$\Lambda_D(u,h)=A(u\oplus h)^{-1}\chi_V(\det_Eh)$; for a pure
multiplier twist it is $\kappa^{-1}\circ\lambda_{H_{\mathrm U}}$.
On the common scalar center the ambient representation in
\eqref{gu:eq:hom-general} has character
$\omega_\tau(z)^{-1}\omega_\sigma(\bar z)^{-1}k_D(z)$.
If $\Pi$ is a theta quotient with $\omega_\Pi=\omega_\sigma$, this
is one by \eqref{gu:eq:theta-central}.

\subsection{The reverse module for fixed general data}
For the module $M_D(\Pi)$ and matching action of
Section~\ref{fun:gu-passage}, the scalar and multiplier twists are
\begin{align}
 \omega_{M_D(\Pi)}(z)&=k_D(z)\omega_\Pi(\bar z)^{-1},\\
 M_D(\Pi\otimes(\eta\circ\lambda_{G'_\theta}))
 &\simeq M_D(\Pi)\otimes(\eta^{-1}\circ\lambda_{G_\theta}),
 \label{gu:asm:multiplier-twist}
\end{align}
where $\eta:F^\times\to\mathbb C^\times$ is smooth.
The first formula follows by taking the matching element
$h=\bar zI_W$ in \eqref{fun:matching-action}; the second follows by
twisting its second tensor factor.

Put $H=G'_\theta$, $H_1=(G'_\theta)^1$, and $H_0=E^\times H_1$.
If $\Pi\simeq\Pi\otimes(\eta\circ\lambda_H)$,
its central character forces $\eta|_{N(E^\times)}=1$. Thus $\eta$
is trivial or $\omega_{E/F}$. Write $\xi=\omega_{E/F}\circ\lambda_H$.
The index-two induction formula and Frobenius reciprocity give
\[
 r=\dim\End_{H_0}(\Pi)
   =1+\dim\Hom_H(\Pi,\Pi\otimes\xi).
\]
Here $r\in\{1,2\}$ is the restriction length in
Lemma~\ref{fun:restriction}. The nontrivial multiplier self-twist
therefore occurs precisely when $r=2$, a case included in the proof.

\begin{corollary}[Fixed general splitting data]\label{gu:app:general-assembly}
For fixed $D$ and $\Pi\in\Irr(G'_\theta)$, a constituent of
$\Pi|_{(G'_\theta)^1}$ has nonzero ordinary lift if and only if
$\Theta_D^{\GU}(\tau)\twoheadrightarrow\Pi$ for some
$\tau\in\Irr(G_\theta)$.
\end{corollary}
\begin{proof}
This is Proposition~\ref{fun:fixed-data} for the fixed data $D$.
\end{proof}
